\exercisesection{6}

¶ 1. Every element z of the p-adic field \(Q_{p}\) (p a prime number) may be written uniquely as \(p^{h}\sum_{k=0}^{\infty}c_{k}p^{k}\) , where \(h\in Z\) , \(c_{0}\neq0\) , \(0\leqslant c_{k}<p\) for all \(k\geqslant0\) (``p-adic expansion'' of z).

\begin{enumerate}
\def\labelenumi{(\alph{enumi})}
\tightlist
\item
  Show that, in order that \(z \in \mathbf{Q}\) , it is necessary and sufficient that there exist two integers \(m \geqslant 0, n \geqslant 1\) such that \(c_{k+n} = c_k\) for all \(k \geqslant m\) . (Observe that, if \(z = a / b \in \mathbf{Q}\) where \(b\) is not divisible by \(p\) , then
\end{enumerate}

\[
a - b \sum_ {k = 0} ^ {n} c _ {k} p ^ {k} = a _ {n + 1} p ^ {n + 1}
\]

where \(a_{n+1}\) is an integer and the sequence \(|a_k|\) (ordinary absolute value) is bounded.)

\begin{enumerate}
\def\labelenumi{(\alph{enumi})}
\setcounter{enumi}{1}
\tightlist
\item
  Suppose that the increasing sequence \((k_n)\) of integers \(k\) such that \(c_k \neq 0\) is such that \(\limsup_{n \to \infty} (k_{n+1}/k_n) = +\infty\) . Show that \(z\) is transcendental over \(\mathbf{Q}\) . (For \(x \in \mathbf{Q}\) , let \(|x|\) and \(|x|_p\) respectively denote the usual absolute value and a \(p\) -adic absolute value. Suppose that \(\mathrm{P} \in \mathbf{Z}[\mathrm{X}]\) is an irreducible polynomial such that \(\mathrm{P}(z) = 0\) ; if \(z_n = p^n \sum_{k=0}^{\infty} c_k p^k\) , show that \(|\mathrm{P}(z_n)/(z - z_n)|_p\) tends to a limit \(\neq 0\) as \(n\) tends to \(+\infty\) , using Taylor's formula; then obtain a contradiction from the hypothesis considering the usual absolute values \(|\mathrm{P}(z_n)|\) .)
\end{enumerate}

\begin{enumerate}
\def\labelenumi{\arabic{enumi}.}
\setcounter{enumi}{1}
\tightlist
\item
  Let D be a non-commutative field of characteristic \(\neq2\) and Z its centre. Suppose that every commutative subfield K of D containing Z is of degree \(\leqslant2\) over Z.
\end{enumerate}

\begin{enumerate}
\def\labelenumi{(\alph{enumi})}
\item
  Show without using Lemma 1 that \([\mathbf{D}:\mathbf{Z}] = 4\) . (Let \(a\in \mathbf{D} - \mathbf{Z}\) be such that \(a^2\in \mathbf{Z}\) and let \(\sigma (x) = axa^{-1}\) for all \(x\in \mathbf{D}\) ; note that \(\mathbf{D}\) decomposes as a direct sum of two vector subspaces over \(\mathbf{Z}\) , \(\mathrm{D}_{+}\) and \(\mathrm{D}_{-}\) such that \(\sigma (x) = x\) on \(\mathrm{D}_{+}\) , \(\sigma (x) = -x\) on \(\mathrm{D}_{-}\) ; note also that \(\mathrm{D}_{+}\) is a subfield of \(\mathbf{D}\) and \(\mathrm{D}_{-}\) is a 1-dimensional vector space over \(\mathrm{D}_{+}\) ; finally, show that \(\mathrm{D}_{+}\) cannot be distinct from \(\mathbf{Z}(a)\) .)
\item
  Show that \(\mathbf{D}\) is a quaternion algebra over \(\mathbf{Z}\) . (Form a basis of \(\mathbf{D}\) over \(\mathbf{Z}\) with the aid of (a).)
\end{enumerate}

